\subsection{同构与结构的迁移}

让 \(\Sigma\) 是理论中的一类结构 \(\mathfrak{C}\) ，在 \(n\) 主基集 \(x_{1},\ldots ,x_{n}\) ，以 \(m\) 辅助基集 \(\mathbf{A}_1,\dots ,\mathbf{A}_m\) . 设 \(\mathbf{S}\) 设为关于 \(n + m\) 字母的阶梯构造方案，该方案出现在以下典型刻画中 \(\Sigma\) ，并且让 \(\mathbf{R}\) 是……的公理 \(\Sigma\) 。在一个理论中 \(\mathfrak{C}'\) 该理论强于 \(\mathfrak{C}\) ，让 \(\mathbf{U}\) 是物种的一个结构 \(\Sigma\) 关于集合 \(\mathbf{E}_1,\ldots ,\mathbf{E}_n\) （作为主基集）并令 \(\mathbf{U}'\) 是集合上同一种类的结构 \(\mathbf{E}_1',\ldots ,\mathbf{E}_n'\) 。最后，令 \(f_{i}\) （在 \(\mathfrak{C}'\) ）是 \(\mathbf{E}_i\) 到 \(\mathbf{E}_i'\) 上的双射 ( \(1 \leqslant i \leqslant n\) )。然后 \((f_1, \ldots, f_n)\) 被称为这些集合之间的同构 \(\mathbf{E}_1, \ldots, \mathbf{E}_n\) ，赋予该结构 \(\mathbf{U}\) 到集合 \(\mathbf{E}_1', \ldots, \mathbf{E}_n'\) ，赋予该结构 \(\mathbf{U}'\) ，如果我们有（在 \(\mathfrak{C}'\) )

\[
\langle f _ {1}, \dots , f _ {n}, \operatorname{Id} _ {1}, \dots , \operatorname{Id} _ {m} \rangle^ {\mathrm{s}} (\mathrm{U}) = \mathrm{U} ^ {\prime}\tag{4}
\]

其中 \(Id_{h}\) 表示 \(A_{h}\) 到自身的恒等映射（ \(1 \leqslant h \leqslant m\) ）。¶ 设 \(f_{i}^{\prime}\) 是双射 \(f_{i}\) 的逆映射（ \(1 \leqslant i \leqslant n\) ）。由（4）和准则 CST3（第2号）立即得到

\[
\langle f _ {1} ^ {\prime}, \dots , f _ {n} ^ {\prime}, \operatorname{Id} _ {1}, \dots , \operatorname{Id} _ {m} \rangle^ {\mathrm{s}} (\mathbf {U} ^ {\prime}) = \mathbf {U}
\]

，因此 \((f_1', \ldots, f_n')\) 是……的一个同构 \(\mathbf{E}_1', \ldots, \mathbf{E}_n'\) ，赋予 \(\mathbf{U}'\) ，映到 \(\mathbf{E}_1, \ldots, \mathbf{E}_n\) ，赋予 \(\mathbf{U}\) 。这些同构 \((f_1, \ldots, f_n)\) 并且 \((f_1', \ldots, f_n')\) 被称为彼此的逆。

\(E_{1}^{\prime}, \ldots, E_{n}^{\prime}\) ，赋予 \(U^{\prime}\) ，若存在 \(E_{1}, \ldots, E_{n}\) 的一个同构，则称其与 \(E_{1}, \ldots, E_{n}\) 到上 \(E_{1}^{\prime}, \ldots, E_{n}^{\prime}\) （赋予U）同构；在这种情况下，结构U与 \(U^{\prime}\) 被称为同构的。上述定义，连同CST1，蕴含以下判据：

CST4. 设 U， \(U'\) , \(U''\) 是同一物种的三个结构 \(\Sigma\) 在主基集上 \(E_{1}, \ldots, E_{n}, E_{1}', \ldots, E_{n}', E_{1}'', \ldots, E_{n}''\) 是物种 \(f_{i}\) 是……的双射 \(E_{i}\) 到上 \(E_{i}'\) ，并且让 \(g_{i}\) 是……的双射 \(E_{i}'\) 到上 \(E_{i}''\) ( \(1 \leqslant i \leqslant n\) ）。如果 \((f_{1}, \ldots, f_{n})\) 并且 \((g_{1}, \ldots, g_{n})\) 都是同构，则 \((g_{1} \circ f_{1}, \ldots, g_{n} \circ f_{n})\) .

的一个同构 \(E_{1}, \ldots, E_{n}\) 到上 \(E_{1}, \ldots, E_{n}\) （关于同一结构）称为 \(E_{1}, \ldots, E_{n}\) 的自同构。 \(E_{1}, \ldots, E_{n}\) 的两个自同构的复合仍是自同构，自同构的逆也是自同构，* 因此 \(E_{1}, \ldots, E_{n}\) 的自同构构成一个群。*

注记。按照语言的滥用，如果 \(f_{i}\) 是...的任意双射 \(\mathbf{E}_i\) 到上 \(\mathbf{E}_i'\) ( \(1 \leqslant i \leqslant n\) ), \((f_1, \ldots, f_n)\) 被称为...的一个同构 \(\mathbf{E}_1, \ldots, \mathbf{E}_n\) 到上 \(\mathbf{E}_1', \ldots, \mathbf{E}_n'\) 关于集合的结构种（第4号，注记3）。

CST5. 在一个理论 \(\mathfrak{G}'\) 该理论强于 \(\mathfrak{G}\) 中，设U是物种 \(\Sigma\) 在 \(\mathbf{E}_1, \ldots, \mathbf{E}_n\) 上的一个结构，并且让 \(f_i\) 是……的双射 \(\mathbf{E}_i\) 到集合 \(\mathbf{E}_i'\) ( \(1 \leqslant i \leqslant n\) ). 那么存在唯一的物种 \(\Sigma\) 在 \(\mathbf{E}_1', \ldots, \mathbf{E}_n'\) 上的结构，使得 \((f_1, \ldots, f_n)\) 是……的一个同构 \(\mathbf{E}_1, \ldots, \mathbf{E}_n\) 到上 \(\mathbf{E}_1', \ldots, \mathbf{E}_n'\) .

对于这个结构，如果它存在，只能是项 \(U'\) 由关系(4)定义；仍需验证该项确实是物种 \(\Sigma\) 的结构，即关系 \(R\{E_{1}', \ldots, E_{n}', U'\}\) 在 \(C'\) . 但这由以下事实推出： \(R\{x_{1}, \ldots, x_{n}, s\}\) 是可迁移的，因为

\(R\{E_{1}^{\prime},\ldots,E_{n}^{\prime},U^{\prime}\}\) 在 \(C'\) 中等价于关系 \(R\{E_{1},\ldots,E_{n},U\}\) （第3号），而该关系在 \(C'\) 中由假设为真。

该结构 \(U'\) 被称为通过将结构U传输到集合而得到的 \(E_{1}^{\prime}, \ldots, E_{n}^{\prime}\) 通过双射映射 \(f_{1}, \ldots, f_{n}\) 因此，同一类型的两个结构是同构的，当且仅当每一个都是通过结构的迁移从另一个得到的。

可能发生这样的情况：任意两个类型为 \(\Sigma\) 的结构必然是同构的；结构类型 \(\Sigma\) 则称为单叶的。* 这就是无限循环群的结构（同构于 \(\mathbf{Z}\) ），特征为零的素域的结构（同构于 \(\mathbf{Q}\) ），完备阿基米德有序域的结构（同构于 \(\mathbf{R}\) ），一个代数闭的、连通的、局部紧致拓扑域的结构（同构于 \(\mathbf{C}\) ),以及非交换、连通、局部紧拓扑除环（同构于 \(\mathbf{H}\) ，即四元数除环）。对于其中某些结构类型，例如特征为零的素域，或完备的阿基米德有序域，甚至不存在除恒等映射以外的任何自同构；但对于上述给出的其他例子，确实存在这样的自同构（例如，对称 \(x \to -x\) 在 \(\mathbf{Z}\) ). *

可以观察到，上述结构种类本质上正是经典数学基础所依据的那些结构。另一方面，* 群结构的种类、有序集结构的种类以及拓扑结构的种类，并不是单价的。*

